\section{Introduction}
\label{sec:introduction}

The rank of a projective module over a group ring contains information
about individual conjugacy classes. This refinement of ordinary rank is
a basic ingredient of group-ring Euler characteristics: alternating sums
of projective-module ranks retain information that a numerical Euler
characteristic forgets \cite{Bass1976,Bass1979}. The Bass trace conjecture
asks how much of that information can occur at elements of infinite order.
We first specify the complex version of this question.

For a unital complex algebra $A$, let $K_0(A)$ be the Grothendieck group of
finitely generated projective right $A$-modules, and set
\[
 \HH_0(A)=A/[A,A],
\]
where $[A,A]$ is the complex linear span of commutators. An idempotent
$e=(e_{ij})\in M_q(A)$ represents the projective right module
$eA^q$, with matrices acting on columns by left multiplication. Its
Hattori--Stallings trace is
\[
 \HS_A([eA^q])=\left[\sum_{i=1}^q e_{ii}\right]\in\HH_0(A),
\]
extended additively to virtual projective classes.

Let $G$ be a discrete group and write $\con(G)$ for its set of conjugacy
classes. The quotient $\HH_0(\C G)$ identifies with
$\bigoplus_{C\in\con(G)}\C C$ by
\begin{equation}\label{eq:class-trace}
 \left[\sum_{h\in G}a_hh\right]
 \longmapsto
 \sum_{C\in\con(G)}\left(\sum_{h\in C}a_h\right)C.
\end{equation}
Indeed, a commutator has zero sum on each conjugacy class, and the
differences of conjugate group elements span the kernel of the displayed
map. Such differences are commutators: $tst^{-1}-s=[t,st^{-1}]$.
All sums arising from an element of $\C G$ have finite support.
We write $\HS_G(x)(C)$ for the coefficient at $C$.

The complex group-ring Bass conjecture asks whether these coefficients
vanish on every infinite-order conjugacy class, for every group and every
class in $K_0(\C G)$; see
\cite[Conjecture 2.1]{BCM2004} and \cite[Conjecture 2]{CM2011}.

\begin{theorem}\label{thm:bass}
Let $G$ be any discrete group. For every $x\in K_0(\C G)$ and every
infinite-order element $g\in G$,
\[
 \HS_G(x)([g]_G)=0.
\]
Equivalently, if $\con_{\mathrm{fin}}(G)$ denotes the conjugacy classes of
finite-order elements, then
\[
 \HS_G\bigl(K_0(\C G)\bigr)
 \subseteq\bigoplus_{C\in\con_{\mathrm{fin}}(G)}\C C.
\]
\end{theorem}

This resolves the complex group-ring formulation positively. Finite-order
coefficients need not vanish away from the identity: for a finite group
$F$, the idempotent $|F|^{-1}\sum_{h\in F}h$ has coefficient
$|C|/|F|$ at every conjugacy class $C$. The conclusion therefore differs
from the identity-only support assertion for integral group rings.
That integral assertion follows in Section~\ref{sec:completion} by combining
Theorem~\ref{thm:bass} with the known restriction on its finite-order
coefficients.

\subsection{Trace rank and the idempotent conjecture}

For a torsion-free group the identity is the only finite-order element,
so Theorem~\ref{thm:bass} also has numerical and ring-theoretic
consequences. Define the canonical trace and the augmentation by
\[
 \tau_G\!\left(\sum_{g\in G}a_g g\right)=a_1,\qquad
 \epsilon_G\!\left(\sum_{g\in G}a_g g\right)=\sum_{g\in G}a_g.
\]
For a matrix $b=(b_{ij})\in M_q(\C G)$, put
$\tau_{G,q}(b)=\sum_{i=1}^q\tau_G(b_{ii})$, and let
$\epsilon_{G,q}(b)\in M_q(\C)$ be the entrywise augmentation.
The coefficient at $[1]_G$ of the Hattori--Stallings trace defines
\[
 \tau_{G,*}:K_0(\C G)\longrightarrow\C,\qquad
 \tau_{G,*}([e(\C G)^q])=\tau_{G,q}(e).
\]
The augmentation also defines the integer-valued augmentation rank
\[
 \operatorname{rk}_{\epsilon}:K_0(\C G)
 \xrightarrow{(\epsilon_G)_*}K_0(\C)
 \xrightarrow{\dim_{\C}}\Z.
\]

The idempotent conjecture commonly attributed to Kaplansky asks whether,
for every torsion-free group $G$ and every field $k$, an element
$e\in kG$ satisfying $e^2=e$ must be $0$ or $1$; see
\cite[Problem 3]{OW2023}. It concerns elements of the group ring itself,
not arbitrary idempotent matrices over that ring.

\begin{corollary}\label{cor:torsion-free}
Let $G$ be any torsion-free discrete group.
\begin{enumerate}[label=\textup{(\roman*)}]
 \item For every $q\geq1$ and every idempotent $e\in M_q(\C G)$,
 \[
  \tau_{G,q}(e)=\operatorname{rank}_{\C}\bigl(\epsilon_{G,q}(e)\bigr).
 \]
 Consequently $\tau_{G,*}=\operatorname{rk}_{\epsilon}$, with
 $\Z$ viewed as a subgroup of $\C$, and
 \[
  \tau_{G,*}\bigl(K_0(\C G)\bigr)=\Z.
 \]
 \item For every commutative unital domain $R$ of characteristic zero,
 every $e\in RG$ satisfying $e^2=e$ is either $0$ or $1$.
\end{enumerate}
\end{corollary}

Part \textup{(ii)} resolves the characteristic-zero Kaplansky idempotent
conjecture positively, and includes coefficient domains as well as fields.
Section~\ref{sec:torsion-free} derives both parts from the trace theorem.
The integer range in part \textup{(i)} concerns algebraic $K_0(\C G)$;
its relation to the reduced $C^*$-completion is made explicit there.

The integral trace consequence also links the theorem to fixed-point
theory. A map $f$ is a \emph{homotopy idempotent} if $f\circ f\simeq f$.
Berrick, Chatterji and Mislin's equivalences \cite[Theorems 1--2]{BCM2006}
then imply that every homotopy idempotent on a connected closed smooth
oriented manifold of dimension greater than two is freely homotopic to
a map with exactly one fixed point. They also identify ordinary and
$L^2$-Lefschetz numbers for pointed homotopy idempotents of connected
closed smooth oriented manifolds that induce the identity on the
fundamental group.
Corollary~\ref{cor:homotopy-idempotents} states these consequences precisely;
Section~\ref{sec:completion} explains their connection to projective
Euler characteristics.

\subsection{The two parts of the proof}

We first describe the algebraic construction to explain the geometric
statement it needs. The body of the paper begins with that geometric
statement, formulated independently of group rings, and completes its
proof before constructing the cycles and cocycles to which it applies.

For a finitely generated group $G$, the algebraic construction starts
with an idempotent $e\in M_q(\C G)$ and one infinite-order element
$g\in G$. The finitely many nonzero group coefficients of $e$ determine
a finite set of steps. Products of matrix
coefficients along paths from $h$ to $gh$, followed by the ordinary
matrix trace, give complex weights on finite lists of group elements.
Idempotency says that summing over any intermediate vertex removes that
vertex. Cyclicity of matrix trace gives a second identity, a cyclic shift
whose last vertex is multiplied by $g$.

Project the vertices to the coset set $Y=\langle g\rangle\backslash G$.
The group $H=C_G(g)/\langle g\rangle$, where $C_G(g)$ is the centralizer
of $g$, acts freely on $Y$. For every positive even dimension $2m$, the weighted
lists define a simplicial cycle modulo this action. Its simplices use
connected sets of vertices in one graph of bounded degree, fixed as $m$
increases. A Pfaffian cocycle evaluates on this cycle as
$2^{-m}\HS_G([e(\C G)^q])([g]_G)$. Here a cocycle means an alternating
function on vertex lists that vanishes on simplicial boundaries. The
explicit construction and the evaluation occupy
Sections~\ref{sec:idempotents} and~\ref{sec:detector}.

The geometric part, proved first in Sections~\ref{sec:framework}--\ref{sec:chains},
shows that every invariant cocycle evaluates to zero on cycles of this
kind in sufficiently high dimension. We cut the supporting simplicial
complex successively into pieces of bounded diameter, choosing the
separating sets nearly to minimize their areas. A volume estimate and
repeated slicing make the expected number of points in the final
zero-dimensional set arbitrarily small. A subdivision construction turns
this point count into a bound for cocycle evaluations. Equivariant chain
homotopy keeps the evaluation unchanged throughout the construction.

The fixed graph degree makes increasing the dimension useful. Through
each vertex there are only exponentially many connected vertex sets of a
given size. On an $n$-simplex, Euclidean structures chosen from spanning
trees of its vertex set give volume $1/n!$; repeated slicing produces the
same factorial in the comparison with a point of the final separating
set. After these factors cancel, radii chosen in terms of the fixed
degree bound make the remaining volume ratio tend to zero with $n$.
The local replacement argument turns this strict volume margin into the
small expected point count. No homological-dimension bound on $H$ is
assumed.

Two features permit this argument for arbitrary groups. A compact
probability space with a free action separates overlapping translates of
local modifications, without any amenability assumption. Also, each top
simplex has its own averaged Euclidean area functional. The areas need
not match along common faces; the successive cuts keep the relevant
dimensional measures away from those faces.
The cocycles need not be globally bounded: near each original simplex,
the construction uses only a fixed finite set of possible vertices.
Bounding the cocycle on those lists suffices for the evaluation estimate.

The final assembly in Section~\ref{sec:completion} chooses a large even
dimension, forcing the individual trace coefficient to vanish. Passing
from finitely generated groups to arbitrary groups uses the exact sum
over subgroup conjugacy classes contained in an ambient conjugacy class.
Section~\ref{sec:torsion-free} then extracts the augmentation rank from
the identity coefficient. In the scalar case this rank is $0$ or $1$,
and the standard faithful-trace argument forces the idempotent itself
to be $0$ or $1$. A finite-coefficient reduction extends that conclusion
from complex coefficients to every commutative characteristic-zero domain.

\subsection{Relation to previous work}

The rank construction comes from Hattori and Stallings
\cite{Hattori1965,Stallings1965}; Bass developed its relation to Euler
characteristics and formulated the integral conjecture in
\cite{Bass1976}; see also \cite[Conjecture (4.5)]{Bass1979}.
Bass proved complex trace vanishing at infinite-order elements for linear
groups \cite[Theorem (6.1)(f)]{Bass1979}. For the integral trace,
Linnell established a restriction on the possible nonzero coefficients
\cite[Lemma 4.1]{Linnell1983}, later recovered through topological
Hochschild homology by Berrick and Hesselholt \cite[Theorem A]{BH2015}.
Its finite-order consequence is the input used in
Section~\ref{sec:completion}.

The scalar idempotent problem also has a long history of trace methods.
Formanek proved the characteristic-zero conclusion for torsion-free
groups satisfying the ascending chain condition on subgroups
\cite[Theorem 1]{Formanek1973}. Burger and Valette revisit the theorems of
Kaplansky and Zalesskii: the canonical trace of a complex scalar
idempotent lies in $[0,1]$, its endpoints force the idempotent to be
trivial, and its value is rational
\cite[Theorems 1.3 and 2.1]{BV1998}. They also obtain
$\tau_{G,*}(K_0(\C G))\subseteq\mathbb Q$
\cite[Corollary 1.2]{BV1998}. Rationality still permits nonintegral
values strictly between zero and one. For torsion-free groups,
Theorem~\ref{thm:bass} identifies the trace with augmentation rank,
which excludes those values. The subsequent faithful-trace argument
is classical and is included in Section~\ref{sec:torsion-free}.

Broader group classes have been treated by algebraic and analytic
methods. Farrell and Linnell proved the complex Bass statement for
elementary amenable groups as part of a result over integral domains
\cite[Theorem 1.6]{FarrellLinnell2003}. Berrick, Chatterji and Mislin
obtained it for amenable groups through rational surjectivity of the
degree-zero Bost assembly map for countable groups
\cite[Theorems 1.3--1.4 and Lemma 1.5]{BCM2004}. Their result concerns
the stronger trace assertion on $\ell^1(G)$, with finite support on
finite-order classes. An algebraic route uses induction from finite
subgroups: the Farrell--Jones conjecture with complex coefficients
implies the complex Bass statement
\cite[Theorems 1.5 and 1.9]{BartelsLuckReich2008}.
The present proof uses the finite support of group-ring idempotents
throughout; it requires neither an assembly hypothesis nor a statement
about a completed group algebra.

For some groups, analytic methods give a stronger idempotent conclusion
in the reduced group $C^*$-algebra $C_r^*(G)$, the operator-norm closure
of $\C G$ acting on $\ell^2(G)$. Using Lafforgue's Banach $KK$-theory,
Mineyev and Yu proved this for torsion-free subgroups of hyperbolic groups
\cite[Theorems 19--21]{MineyevYu2002}. Puschnigg independently obtained
the torsion-free hyperbolic case through local cyclic homology
\cite{Puschnigg2002}; see also \cite[Section 6]{MineyevYu2002}.
Our conclusion for arbitrary torsion-free groups concerns the group ring
and the trace on its algebraic $K_0$, as distinguished in
Section~\ref{sec:torsion-free}.

The higher-dimensional algebraic construction has antecedents in the
cyclic-homological approach to the conjecture. Connes constructed a pairing
of even cyclic cocycles with projective classes
\cite[Part II, \S2, Proposition 14]{Connes1985}. Burghelea's decomposition
of cyclic homology by conjugacy classes identifies the infinite-order
components through the reduced centralizers $C_G(g)/\langle g\rangle$
\cite[Proposition 1.8 and Theorem I]{Burghelea1985}. Eckmann used
higher-degree vanishing to obtain trace vanishing over rational group
rings for specified classes of groups of finite rational homological
dimension \cite[\S3, Theorems 3.1 and 3.3]{Eckmann1986}.
Sections~\ref{sec:idempotents}--\ref{sec:detector} give direct finite-chain
constructions in this tradition, including the normalization of the
Pfaffian detector; no cyclic-homology theorem is needed as an input.
The sparse-chain theorem replaces a homological-dimension hypothesis
with connected-support control on the particular cycles.

The geometric proof develops the separating-set method of Papasoglu
\cite[\S2, especially Lemma 2.5]{Papasoglu2020}, the separating filtrations and simplex-counting
argument of Alpert \cite[Lemmas 4, 6--7]{Alpert2025}, and the
Cantor-parameter approach of Braun and Sauer
\cite[\S\S2, 4--5]{BraunSauer2021}. Alpert's equivariant formulations
\cite[\S4, especially Lemmas 11--12]{Alpert2025} provide the closest model
for the local comparisons. Alpert's filtration results concern manifold
geometry and compatible metrics. We give the full piecewise-linear argument needed
for simplex-dependent areas and arbitrary cycles in the coinvariant
chain complex. The separate Pfaffian calculation converts this
geometric vanishing into a statement about each conjugacy-class
coefficient.
